\documentclass{article}
\usepackage{graphicx}

\title{TDSE Project}

\author{
  Juan Pablo Vélez Muñoz
  \and
  Luiza Mariana Gonzalez Veloza
  \and
  Daniel Alexander Ahumada León
}

\date{August 2026}

\begin{document}

\maketitle

% --- INTRODUCTION ---
\section{Introduction}

Keeping track of environmental and social issues in remote areas is hard because reaching these places is difficult. Satellite images help us observe large areas of land over time, but checking all these images manually takes too much time and effort.

In this project, we look at how to use automated tools to analyze satellite images. The goal is to spot interesting patterns and highlight areas that might need a closer look, helping teams focus on where it matters most.

Here is how we organized this paper: we explain the main problem, review what others have built, present our proposed system, and show our plan for the prototype and testing.


\subsection{Context}
Monitoring large territories takes a lot of time and resources. While satellite images give us a way to look at changes from above, handling thousands of images manually creates a bottleneck. We need a simpler, faster way to process this data.

\subsection{Problem / Niche}
Environmental and social activities can quickly change large areas of land, especially in remote places where sending people to check on-site is hard and expensive. Satellite images cover a lot of ground, but there are so many of them that checking everything manually is almost impossible.

Our focus is on using automated tools to scan satellite images and highlight areas where important changes are happening. The main challenge is making sure the system gives useful alerts without generating too many false alarms, and without assuming every change on the map means something illegal is going on.
\subsection{Contributions}
% What we are building and bringing to the table

\subsection{Document Structure}
% How the paper is organized


% --- BODY OF THE PAPER ---
\section{Problem}
\subsection{Qualitative Aspects}
\subsection{Quantitative Aspects}

\section{State of the Art}
\subsection{Grey Literature}
\subsection{Scientific Literature}
\subsection{Technical Documentation}

\section{Motivation / GAP / Solution}

\section{Architecture}

\section{Prototype}

\section{Results}
\subsection{Qualitative Results}
\subsection{Quantitative Results}


% --- CONCLUSIONS AND FUTURE WORK ---
\section{Conclusions and Future Work}
\subsection{Contributions}
\subsection{Final Context}
\subsection{Future Work}


\bibliographystyle{plain}

\end{document}